\section{Geometry and the two distinguished quotients}
We use the classical marked Coxeter group $[5,3]$ of order $120$ for the dodecahedron, its reversed marking for the icosahedron, and the antipodal quotient of order $60$ for the corresponding hemi-polyhedron. The central antipode is the fifth power of the product of the three distinguished reflections. We also use the standard string-C-group/polytope correspondence and the general semisparse quotient correspondence. These are retained geometric inputs, not the target universal-group theorem.

\begin{theorem}\label{geometry}
The group in Theorem~\ref{structure}, with its displayed marking, is a string C-group. Its regular polytope $\mathcal P$ has dodecahedral facets and hemi-icosahedral vertex figures. It is the universal polytope with these local types. The quotient with group $J_1$ has $1463$ dodecahedral facets and $2926$ vertices; the quotient with group $\PSL_2(19)$ is the self-dual $57$-cell with hemi-dodecahedral facets and hemi-icosahedral vertex figures. Both distinguished quotients have no proper polytope quotient under the semisparse convention.
\end{theorem}
\begin{proof}
For each of the three actual markings, the required program constructs the subgroup on every subset of the four markers and checks all $16^2=256$ intersection equalities. The markers are nonidentity involutions, their adjacent products have exact orders $5,3,5$, and the nonadjacent pairs commute. These checks establish the entire string C-property, not only its two endpoint intersections.

The endpoint orders are respectively $(120,60)$ for $W$ and $J_1$, and $(60,60)$ for the projective quotient, where the first entry is the facet group. Their Coxeter presentations give surjections from the classical local groups just described, and equality of their orders makes them marked isomorphisms. Thus their local polytopes are the claimed ones; an abstract group isomorphism to $C_2\times A_5$ alone would not suffice. The string-C-group correspondence constructs the regular polytopes. Group order divided by the appropriate endpoint order gives $1463$, $2926$ and $57$. Reversing the four actual projective markers extends to a bijective group homomorphism, explicitly checked by the program, and hence proves the asserted self-duality.

For universality, take any polytope with the indicated local types and let its involutory flag adjacencies be $r_i$. Nonadjacent adjacencies commute; the polygon sections give the three adjacent relations. At each vertex the hemi-icosahedral flag action gives $(r_1r_2r_3)^5=1$. Thus \eqref{W} acts transitively on the flags, and the stabilizer of a chosen flag identifies the polytope with the corresponding quotient of $\mathcal P$ by the general quotient correspondence. This applies without assuming the target regular. It also applies when some facets have been replaced by their antipodal quotients, since their adjacencies still satisfy the same facet relations. Every section-regular locally projective polytope of type $\{5,3,5\}$ is of one of these local types or has dual of one of them. This proves the precise classification coverage claimed in the source; the mixed-facet quotients below are additional, non-section-regular members of the full quotient census.

For the final assertion retain the classical fact that the hemi-icosahedron has no proper quotient and Hartley's general endpoint semisparseness theorem. Applied to these local groups, it says that a subgroup $N$ of the Janko marked group is semisparse exactly when
\[
 N^g\cap H_0H_3\subseteq\{1,\omega\}\quad\hbox{for all }g;
\]
in the projective quotient the right side is $\{1\}$. The facet alternatives are respectively the dodecahedron and its antipodal quotient, or only the hemi-dodecahedron. For each prime and every conjugacy class of prime-order subgroups, the finite program finds a member meeting the corresponding forbidden endpoint product. Lemma~\ref{prime} therefore excludes every nontrivial semisparse subgroup.

This calculation also checks the source's explicit Janko witnesses: in its marker indices the words
\[
 s_1s_2,\quad s_0s_1,\quad s_3s_1s_2s_1s_0,\quad
 s_3s_2s_1s_0s_1s_0,\quad
 s_1s_3s_2s_1s_3s_0s_1s_0s_2s_1s_0
\]
belong to $H_0H_3$ and have orders $3,5,19,11,7$. Its unique involution class and its prime-order odd Sylow groups give the original short Cauchy argument as well. For the projective quotient the orders tested are $2,3,5,19$; the test is on prime-order subgroups, so does not incorrectly require a Sylow $3$-subgroup to have order three.
\end{proof}

